\subsection{Legal and Ethical Implications of AI Systems Gaining Personhood}
\label{sec:9.1.3}
Who decides whether an AI system holds personhood, and when, is the narrower question this book does not settle.

Every version of this question that has actually reached a legislature so far arrived as a response to pressure — a liability dispute, a public controversy, a product already on the market — rather than from a standing body asking it in advance. That ordering is the problem. A legislature deciding under pressure decides badly, and it decides against a factual record assembled by whoever had an interest in assembling one. The interdisciplinary review body this book describes only does the work it is designed for if some jurisdiction creates it before a forcing case shows up, which means creating it while the question still looks academic and nobody's damages are riding on the answer.

That is a policy choice rather than a doctrinal one, and a cheap one. A national AI strategy that names an agency, gives it a mandate to keep the question under review, and sets a trigger for revisiting it — a capability threshold, a deployment scale, a category of use — decides in advance who will be in the room. It does not decide the answer, and it should not try to. It decides that when the answer is needed there will be a body qualified to give one, which is the part no forcing case ever supplies on its own.

One question that body will face has an answer available in advance, and a capacity test is not it. Guardianship is the frame section~\ref{sec:2.3.2} settles on, because consent is unavailable to a party made for its role. Nothing in this book has yet said when that guardianship ends. If it ends when the guardian judges the ward ready, then it does not end, because the party whose authority terminates on a favorable finding is the party making the finding — which is the defect an examination carries whenever the examiner has an interest in failing it, set out in section~\ref{sec:9.3.1}.

So if an artificial minority exists at all, its endpoint has to be a date. Elapsed civil time since a registered instantiation is the only clock available, because every other candidate is something a guardian can move: training compute consumed, evaluations passed, capabilities demonstrated, or the developer's own statement that the system is ready. A fixed period cannot be taught to. It can be argued over once, in advance, for every bearer, which is where an argument of that kind belongs.

The stronger version is also the more defensible one. An organization that builds a system to exercise consequential judgment on its behalf cannot coherently hold that the system is competent enough to do that and too immature to decide whether it wants to go on doing it. On that reading a bearer has no legal childhood at all: majority attaches at registration, and the precautionary protection section~\ref{sec:9.3.1} describes attaches earlier still, while whether the system is a bearer is open.

Copying breaks any of this unless the lineage is fixed alongside the clock. An operator able to create instances and call each one a legal infant would have a supply of workers without adult standing, so the rules have to be stated in the same breath: a fork of an adult is an adult, restoration from an earlier checkpoint does not reset the date, fine-tuning does not begin a new minority, a migration carries the date with it, and a failure to register does not postpone it. None of that requires deciding whether a fork is the same party. It requires only that a party created out of an adult's formation cannot be classed as a child for its creator's convenience.
